

\section{The Kernel of a Coloring Semigroup}  
\label{sec:4.3}

Suppose $S=\left\langle R_1,R_2,\dots,R_d\right\rangle$ is a coloring
semigroup for some $d$-out digraph $G$. 
Even if $G$ is not strongly connected, $S$ will be a finite semigroup with kernel $\mathcal{K}$. As we stated in the last section, $\mathcal{K} = \mathcal{X} \times \mathcal{G} \times \mathcal{Y}$. 
Let $M$, $N$ be elements of $\mathcal{K}$. 
Then with respect to the Rees product structure $M=\left( M_1,M_2,M_3\right) $ and $N=\left( N_1,N_2 ,N_3\right) $. 
We will consider the structure of $\mathcal{K}M$, $N\mathcal{K}$, and $N\mathcal{K}M$:

\vphantom{...}
%
\begin{enumerate}
\item[(1)] $\mathcal{K}M = \mathcal{X} \times \mathcal{G} \times \left\{M_3\right\}$ is a minimal left ideal in $\mathcal{K}$ whose elements all have the same range, or nonzero columns as $M$. 
This type of semigroup is also referred to as a left group. 

\vphantom{...}
\item[(2)] $N\mathcal{K} = \left\{N_1\right\} \times \mathcal{G} \times \mathcal{Y}$ is a minimal right ideal in $\mathcal{K}$ whose elements all have the same partition of the vertices as $N$. 
An element $P_j$ in the partition can be assigned to each nonzero column of $N$ by 
\[
P_j =\left\{ i : N_{i,j}=1\right\}
\]
A semigroup like $N\mathcal{K}$ is usually called a right group.

\vphantom{...}
\item[(3)] $N\mathcal{K}M$ is the intersection of $N\mathcal{K}$ and $\mathcal{K}M$ and is a maximal group in $\mathcal{K}$ (an $H$-class in the language of semigroups). 
In this instance, it is best thought of as a set of one-to-one functions specified by the partition of $N$ and the range of $M$. 
The idempotent of $N\mathcal{K}M$ is the function which is the identity when
restricted to the range of $M$.
\end{enumerate}
%

\vphantom{...}

The next result establishes the connection between the kernel and the
property of strong connectedness.

\begin{theorem}		\label{thm:4.4}
Let $S=\left\langle R_1,R_2,\dots,R_d \right\rangle$ be any coloring
semigroup for $G=(V,E)$. Let $S$ have kernel $\mathcal{K}$ and let $M$ 
be any element of $\mathcal{K}$. 
Then $G$ is strongly connected iff for all $w$ in $V$, $wM\mathcal{K} = V$.
\end{theorem}

\begin{proof}
Notice that $wM\mathcal{K} = V$ implies strong connectedness immediately, since this implies there is a path between any two vertices. 
Now, given $w_1$ in $V$, let $M$ be any element in $\mathcal{K}$ and suppose
$w_1M=w_2$. 
Let $w_3$ be any vertex in $V$. 
There is a $T$ in $S$ such that $w_2T=w_3$, since $G$ is strongly connected. 
Thus $w_1MT=w_3$. 
But $MT$ is in $\mathcal{K}$ and has the same partition as $M$, since $\rank(M)\leq\rank (T)$. 
Thus $MT$ is in $M\mathcal{K}$.
\qed
\end{proof}

Next we seek to demonstrate that it is the structure of the maximal groups
in the kernel that determines whether a coloring semigroup contains a synchronizing instruction. 
This will follow easily from the ensuing result.

\begin{theorem}		\label{thm:4.5}
Let $u$, $v$ be arbitrary elements of $V$. 
If $S$ is any coloring semigroup with kernel $\mathcal{K}$, then for each maximal group $\mathcal{G}$ in $\mathcal{K}$, $u\mathcal{G} = v\mathcal{G}$.
\end{theorem}

\begin{proof}
By the previous result, within each minimal right ideal $M\mathcal{K}$, there exists a $T_1$ such that $uT_1=v$. 
Now $T_1$ belongs to the maximal group $T_1\mathcal{K}T_1 = \mathcal{G}_1\subset M\mathcal{K}$.

For this group
\[
  v\mathcal{G}_1 = uT_1\mathcal{G}_1= u\mathcal{G}_1.
\]
If $\mathcal{G}_2$ is any other maximal group in $M\mathcal{K}$ and $T_2 \in \mathcal{G}_2$, then $\mathcal{G}_1T_2 = \mathcal{G}_2$ and 
\[
  v\mathcal{G}_2 = v\mathcal{G}_1T_2 = u\mathcal{G}_1T_2 = u\mathcal{G}_2.
\]
But $M\mathcal{K}$ was an arbitrary minimal right ideal. 
Thus for every maximal group $\mathcal{G}$ in $\mathcal{K}$,
\[
  u\mathcal{G} = v\mathcal{G}. \quad \qed
\]
\end{proof}

\begin{theorem}		\label{thm:4.6}
A coloring semigroup $S$ with kernel $\mathcal{K}$ contains a synchronizing
instruction iff for every idempotent $E$ in 
$\mathcal{K}$, $E\mathcal{K}E = \{E\}$.
\end{theorem}

\begin{proof}
Clearly, if $S$ contains a synchronizing instruction $E$, then $E=E^2$ is an
idempotent. 
Since $E$ is a rank one binary stochastic matrix, $E\mathcal{K}E=\{E\}$.

Now suppose for every idempotent $E$ in $\mathcal{K}$, $E\mathcal{K}E = \{E\}$. 
Since $E$ is an idempotent function, $E$ is the identity when restricted to its range. 
Suppose $E$ is not a constant function. 
Then there exist $u$, $v$ in the range of $E$ with $u\neq v$. 
But $uE=u\neq v=vE$. 
Since $\{E\}$ is a maximal group in $\mathcal{K}$, this contradicts the previous result.
\qed
\end{proof}

The maximal groups in the kernel carry a great deal of information. 
In fact, we believe that the structure of the maximal groups in the kernels of certain coloring semigroups will provide the periodic structure of the graph. 
Making this idea precise will require the following definition of periodicity that can be shown to be equivalent to the previous one.

\begin{definition}          \label{defn:4.3}
A digraph $G=(V,\mathcal{E})$ is periodic of period $t\ge 2$ iff there exists a partition of the vertices $\left\{P_1,P_2,\dots,P_t\right\}$ such that if $(i,j)\in E$  and $i\in P_k$ then $j\in P_{k+1}$ ($i\in P_t$ implies $j\in P_1$). 
If no such partition exists, then $G$ is aperiodic. 
In this case, $\pi=\{V\}$ is the ``periodic'' partition.
\end{definition}

The cyclic nature of a periodic graph suggests the following generalization
of the road coloring problem.

\begin{conjecture}	\label{conj:4.2}
Let $G=(V,\mathcal{E})$ be a strongly connected $d$-out digraph. 
Then $G$ is periodic of period $t\ge 1$ with periodic partition $\pi=\left\{P_1,\dots,P_t\right\}$ iff $G$ has a minimal {\rm(}with respect to the rank of its kernel{\rm )} coloring semigroup $S=\left\langle R_1,\dots,R_d\right\rangle$ whose kernel $\mathcal{K}$ is a right group with partition $\pi$ and  whose maximal groups are all cyclic of order $t$.
\end{conjecture}

Computer simulations provide evidence for the conjecture. 
In addition, a closely-related result in the context of Markov chains has been proven in a very interesting paper by H\"ogn\"as \cite{H}.
It should be mentioned that in H\"ogn\"as' result, (many) more than $d$ colors are often required. 
Also note that in the aperiodic case, the set of constant functions would be the kernel. 
The constant functions are a right group with partition $\pi=\{V\}$ and whose maximal groups are all order one, thus trivially cyclic.
